% ===================================================================== % sections/mechanisms.tex — shared (long + short) % =====================================================================

\section{Mechanisms} \label{sec:mechanisms}

The reduced-form estimates in Section~\ref{sec:results} tell a story about whether early-life exposure to a village midwife moves adult outcomes; they are silent about the pathway through which the program operates. Because the midwife bundle is a composite of prenatal counseling, skilled attendance at delivery, postnatal follow-up, and micronutrient distribution, the reduced-form coefficient conflates every channel into one number. The existing literature has made progress on isolating individual channels: \textcite{frankenbergthomas2001} document post-program increases in maternal body-mass index, \textcite{frankenberg2005} tie the program's arrival to improvements in height-for-age among children born after a midwife was available in the community, and \textcite{ahsan2021} report a positive first-stage effect on village-midwife-attended births and exclusive breastfeeding using IFLS book-4 pregnancy-history data. This section follows the Ahsan first-stage design and asks whether \emph{in-utero} program exposure shifts five mother-reported intermediate outcomes recorded in the IFLS book-4 pregnancy-history module: (i) any prenatal check-up during pregnancy; (ii) any skilled attendant at delivery (a physician, private midwife, village midwife, or nurse); (iii) a \emph{village midwife} specifically as the delivery attendant, the program-specific channel; (iv) exclusive breastfeeding for at least six months; and (v) recorded birth weight in kilograms. The IFLS-4 codebook key for the multi-response delivery-attendant variable separately identifies private midwives (B) from village midwives (C); I report both an aggregate ``any skilled'' indicator and the village-midwife-specific indicator that pins down the program-specific channel, rather than collapsing them into one. The first-stage uses an in-utero exposure indicator equal to one when the village midwife arrived in the community of birth at least one year before the child's birth; this is the cleanest exposure margin for at-birth mediators because children whose program arrived after their birth cannot have these outcomes moved by the program. The sample is restricted to children born in or after 1989, the year the program was launched.

Table~\ref{tab:mechanism_first_stage} reports the first-stage regressions of each mediator on the in-utero-exposure indicator, absorbing community-of-birth, birth-year, and survey-wave fixed effects, with the same mother-level controls used in the main specification. The point estimate on prenatal check-up is essentially zero (\num{-0.004}, SE \num{0.020}), unsurprising given the near-ceiling sample mean of \num{0.936}. The aggregate skilled-attendant indicator is also a precisely-estimated null (\num{0.001}, SE \num{0.032}). The village-midwife-specific indicator, by contrast, picks up a marginally significant first stage of \num{0.045} (SE \num{0.026}, $p \approx 0.08$): an in-utero-exposed child is roughly $4.5$ percentage points more likely to be delivered by a village midwife than an unexposed child in the same community-of-birth and cohort, on a baseline mean of $8.3$ percent. The aggregate skilled and village-midwife rows together imply that the program shifted the composition of skilled deliveries toward village midwives without expanding the overall rate of skilled attendance. The exclusive-breastfeeding-for-six-months indicator and the birth-weight outcome do not move significantly (\num{0.009}, SE \num{0.026}; and $-\num{0.067}$~kg, SE \num{0.053}, respectively), consistent with the village-midwife channel operating through delivery-attendant choice rather than through breastfeeding behavior or anthropometric birth outcomes in this cohort.

\begin{table}[!htbp] \centering \caption{First-stage: effect of in-utero village-midwife exposure on candidate mediators.} \label{tab:mechanism_first_stage} % Auto-generated by code/R/14_mechanisms.R
\begin{tabular}{lcccc}
\toprule
Mediator & Sample mean & $\hat{\beta}$ & SE & $N$ \\
\midrule
Any prenatal check-up & 0.936 & -0.015 & (0.026) & 2,467 \\
Skilled birth attendant & 0.634 & 0.069 & (0.052) & 2,451 \\
Ever breastfed & 0.980 & -0.027$^{*}$ & (0.016) & 2,441 \\
\bottomrule
\end{tabular}
 \begin{minipage}{0.95\textwidth} \footnotesize \emph{Notes.} Each row is a separate TWFE regression of the mediator on the in-utero village-midwife exposure indicator (one if the SAR-recorded rollout year was at least one year before the child's birth year), with mother-level pre-determined controls (sex, mother's years of schooling, mother's age at birth) and community-of-birth, birth-year, and survey-wave fixed effects. Sample is restricted to children born in or after 1989. Mediators are extracted from the IFLS book-4 pregnancy-history module (W2 for 1984--1997 births, W3 for 1998--1999 births); the village-midwife indicator follows the IFLS-4 codebook key in which the delivery-attendant letter code C identifies village midwives specifically. Standard errors clustered at the community-of-birth level in parentheses. Significance: $^{*}\,p<0.10$, $^{**}\,p<0.05$, $^{***}\,p<0.01$. \end{minipage} \end{table}

The first-stage table narrows the channel set to one program-specific margin. The village midwife is the only attendant type whose share moves with in-utero program exposure, on a sample where the aggregate skilled-attendance rate is essentially unchanged; the program's at-the-margin contribution is therefore a substitution \emph{into} village-midwife attendance from other delivery attendants (most plausibly traditional birth attendants and private midwives) rather than a substitution into skilled care from no skilled care. A more ambitious mediation decomposition is not pursued here because the pregnancy-history sub-sample is non-randomly selected and the reduced-form effect on the sub-sample is small and imprecisely estimated, so any share-mediated quantity inherits the imprecision of a small reduced-form numerator and is not a credible channel-share inference. Mother-reported retrospective data on pregnancies ten to thirteen years before the interview also carries recall bias that may attenuate the first stage for precisely the cohorts on whose newly-treated communities the program's effect is identified \autocite{imai2010}.
